% vi: ts=2 sw=2 sts=2 et filetype=tex
% SPDX-License-Identifier: CC-BY-SA-4.0

\documentclass[25pt, a0paper, portrait]{tsj-poster}

% --- Configuración del Estilo ---
\usetheme{TSJ}

% --- Información del Título ---
\title{Sistema de generación y calificación automática de reactivos
de las materias de matemáticas a nivel medio superior}
\author{Yolanda Alejo Huerta , M. Díaz-Rodríguez}
\institute{Tecnológico Superior de Jalisco, Unidad Académica Zapopan,
za240113432@zapopan.tecmm.edu.mx,
miriam.diaz@zapopan.tecmm.edu.mx
}

\begin{document}

\maketitle

\block{Resumen}
{
Este trabajo presenta el desarrollo de un asistente virtual enfocado
en la evaluación automática de competencias técnicas en áreas de
matemáticas a nivel media superior. La propuesta surge ante la
necesidad de optimizar los procesos de enseñanza-aprendizaje en
contextos donde la atención personalizada resulta limitada.
\vspace{1em}

La plataforma, denominada \textbf{Matefácil} versión \textbf{2.0},
implementa una arquitectura basada en \textbf{MVC}
(\emph{Modelo Vista Controlador}), utilizando
tecnologías como \emph{PHP}, \emph{Laravel}, \emph{SQLite}.
Su principal innovación radica en la generación
dinámica de ejercicios matemáticos, específicamente fórmulas
proposicionales, evitando el uso de bancos estáticos de reactivos.
Y su uso en aulas, \emph{sin necesidad de Internet.}
\vspace{1em}

El sistema crea una \textbf{red local/interna} donde el alumno se conecta y
realiza los ejercicios desde \emph{su computadora} o su \emph{telefono móvil}.
Además permite configurar parámetros de complejidad, generando
ejercicios únicos para cada usuario. Adicionalmente integra un módulo de
evaluación automática de estos ejercicios, proporcionando
retroalimentación inmediata.
\vspace{1em}

Los resultados evidencian mejoras en la eficiencia docente y en el
\emph{aprendizaje autónomo} de los estudiantes, posicionando esta
herramienta como una solución escalable para entornos educativos digitales.

\vspace{2em}
{
\color{titlefgcolor}
\bfseries \Large Palabras Clave \par
}
Evaluación automática; matemáticas; generación dinámica; educación digital
}

\begin{columns}
\column{0.5}
\block{I. Introducción}{
En la actualidad, la educación enfrenta el reto de adaptarse a
entornos digitales cada vez más dinámicos, donde el uso de
Tecnologías de la Información y la Comunicación (TIC) es
fundamental para mejorar los procesos de enseñanza-aprendizaje.
\vspace{1em}

Plataformas como Google Classroom, Moodle y los MOOC han permitido
ampliar el acceso a recursos educativos; sin embargo, aún existen
desafíos importantes en la personalización del aprendizaje y la
evaluación continua.
\vspace{1em}

En áreas como la lógica, las matemáticas y la física, los estudiantes
requieren práctica constante mediante la resolución de ejercicios.
No obstante, la generación, evaluación y retroalimentación de estos
ejercicios representa una carga significativa para el docente.
\vspace{1em}

En este contexto, se propone el desarrollo de una plataforma
inteligente que automatiza la generación de ejercicios y su evaluación,
permitiendo ofrecer una experiencia personalizada al estudiante
y optimizando el tiempo del profesor.
}
\block{II. Objetivos}{
{
 \color{titlefgcolor}
 \bfseries \large Objetivo general \par
}

\vspace{1em}
Desarrollar un asistente virtual que permita la generación
y evaluación automática de ejercicios matemáticos para apoyar
la adquisición de competencias técnicas en estudiantes de
áreas computacionales.

\vspace{1em}
{
\color{titlefgcolor}
\bfseries \large Objetivos específicos \par
}

\begin{itemize}
  \item Diseñar una arquitectura web escalable basada en el modelo MTV.
  \item Implementar algoritmos para la generación automática de
        fórmulas proposicionales.
  \item Desarrollar un módulo de evaluación automática basado en
        tablas de verdad.
  \item Permitir la configuración de la complejidad de los ejercicios.
  \item Evaluar el desempeño del sistema en entornos educativos reales.
\end{itemize}

}
    
 \column{0.5}
  \block{Something else}{Here, bla bla bla \vspace{4cm}}
    \note[
        targetoffsetx=-9cm, 
        targetoffsety=-6.5cm, 
        width=0.5\linewidth
        ]
        {e-mail \texttt{welcome@overleaf.com}}
\end{columns}

\end{document}

